% ---------------------------------------------------------------------
%  TÓM TẮT
% ---------------------------------------------------------------------
\begin{abstract}

Ứng dụng mạng tương tác thời gian thực nhiều người dùng là một bài toán cốt
lõi trong kỹ thuật phần mềm hiện đại. Bài tập lớn này xây dựng hệ thống
\term{Trò chơi đoán số nhiều người qua mạng TCP} (Đề tài số 30) trên nền tảng
Java 21. Thách thức trung tâm của bài toán nằm ở việc đồng bộ hóa trạng thái
khi nhiều máy khách cùng gửi lượt đoán đồng thời tới máy chủ, loại bỏ hiện
tượng tranh chấp dữ liệu \en{(race condition)}, phòng tránh bế tắc luồng
\en{(deadlock)} khi người chơi chuyển giữa nhiều phòng chơi song song, và bảo
đảm độ bền vững của máy chủ khi máy khách ngắt kết nối đột ngột giữa ván.

Để giải quyết trọn vẹn các thách thức trên, nhóm thiết kế hệ thống theo kiến
trúc 3 tầng phân tách độc lập gồm Tầng Giao diện và Mạng (\code{ui},
\code{net}), Tầng Nghiệp vụ (\code{domain}, \code{service}) và Tầng Lưu trữ
(\code{repo}). Hệ thống vận dụng mô hình máy chủ đa luồng
\en{(thread-per-connection)} với bể luồng cố định \code{ExecutorService},
giao thức văn bản hướng dòng qua \code{BufferedReader} và \code{PrintWriter}
có bật \code{auto-flush}, cấu trúc bản ghi bất biến \code{record LuotDoan},
kiểu kết quả đóng kín \code{sealed interface KetQua<T>}, hàng đợi hai đầu
\code{ArrayDeque} cho chế độ đoán lần lượt theo vòng, tác vụ hẹn giờ
\code{Future} cho giới hạn thời gian mỗi ván, và quy tắc giành khóa theo thứ
tự từ điển cố định \en{(lock ordering rule)} để chống bế tắc. Toàn bộ 21 chức
năng từ \code{F01} đến \code{F21} đã được phát triển theo quy trình hướng
kiểm thử \en{(Test-Driven Development -- TDD)} và vượt qua \textbf{103/103 ca
kiểm thử tự động JUnit 5} (đạt tỷ lệ 100\%) trong thời gian dưới 6 giây.

\vspace{0.8em}
\noindent\textbf{Từ khoá:} lập trình mạng Socket TCP, Java 21, đồng bộ hóa đa
luồng, chống bế tắc Lock Ordering, Streams API, phát triển hướng kiểm thử TDD.

\end{abstract}
